\documentclass[aspectratio=169]{beamer}

\usetheme{Madrid}
\usepackage{booktabs}
\usepackage{xcolor}
\setbeamertemplate{navigation symbols}{}

\title[IndiaSAT models, biomass, robust water]{Plugging in the lab's production models, biomass, robust water, and non-linear learners}
\author[Salil Gujar]{Salil Gujar \\ \footnotesize M.Tech CSE \quad Entry No.\ 2025MCS2106}
\date{July 2026}

\begin{document}

\frame{\titlepage}

% ------------------------------------------------------------------
\begin{frame}{Where we are, and what this tackled}
\begin{columns}[T]
\begin{column}{0.5\textwidth}
\emph{Done so far}
\begin{itemize}\small
  \item A 4-class base map of India at 10\,m with a living hierarchy: split, add, rule-split, merge,
  retrain a node, all as crisp Earth-Engine tiles.
  \item A git-backed zoo of model and dataset cards, save and resume, GeoTIFF export, per-fortnight
  water.
\end{itemize}
\end{column}
\begin{column}{0.5\textwidth}
\emph{This round}
\begin{itemize}\small
  \item Plug the lab's IndiaSAT models into the framework as usable, carded classifiers.
  \item Add non-linear learners: Random Forest on Alpha Earth, XGBoost on Tessera.
  \item Understand the biomass data collection, and what it would take to use it.
  \item Segment the mining class; make water spatially and temporally robust.
  \item Implement STACD provenance for every output; measure Tessera against Alpha Earth.
\end{itemize}
\end{column}
\end{columns}
\end{frame}

% ------------------------------------------------------------------
\begin{frame}{Plugging in the lab's production models}
Use the IndiaSAT models shared with us: train and store them, and list them in the zoo so a user can
pick one to split a class.
\begin{itemize}\small
  \item Two models: a pan-India \emph{tree against crop} classifier, and a per-region \emph{farm,
  plantation, scrubland} classifier.
  \item We found both train a Random Forest \emph{and} classify entirely inside Earth Engine, so they
  render server-side as tiles, with nothing downloaded and no re-implementation on our side.
  \item Their training data is readable from our project, so we reproduce them faithfully (trained on
  the fly from the stored asset) rather than copying weights.
  \item Each is now a card in the zoo, applied from there.
\end{itemize}
\end{frame}

% ------------------------------------------------------------------
\begin{frame}{The two models, and how they plug in}
\begin{itemize}\small
  \item Tree vs crop, pan-India: features are a Sentinel-1 radar time series (23 sixteen-day steps,
  two polarisations), trained on about seventy thousand labelled points.
  \item Farm / plantation / scrubland, per agro-ecological region: features are the same Alpha Earth
  embedding we already use, trained per region on about 1.5 million points.
  \item Both are a new card topology, \texttt{ee\_rf}, storing the recipe (asset, features, classes),
  not a file, since the model is re-trained in Earth Engine on demand.
  \item They carry their own class scheme, so applying one \emph{rewrites the hierarchy}: greenery
  becomes the model's classes and the rest of the base map (built-up, water, barren) is kept. The tree
  and the map both follow the model.
\end{itemize}
\end{frame}

% ------------------------------------------------------------------
\begin{frame}{Non-linear learners, and where they run}
\begin{itemize}
  \item The aim was Random Forest on Earth Engine, the model that usually works, and XGBoost on
  Tessera. XGBoost on Tessera was almost free; Random Forest on Alpha Earth was the real work.
  \item A Random Forest is not band math, so inference is now \emph{algorithm-aware}: a non-linear
  Alpha Earth split falls back to the point-grid render (the path Tessera already uses), while linear
  models keep the crisp tile map.
  \item This also unlocks biomass, a Random Forest regressor on the same Alpha Earth features.
\end{itemize}
\end{frame}

% ------------------------------------------------------------------
\begin{frame}{Biomass from GEDI: the data collection, and findings}
A study of the biomass scripts we were given, to understand the data collection and how it would fit.
It is not surfaced in the interface.
\begin{itemize}\small
  \item The collection samples GEDI, a lidar that measures above-ground biomass, and pairs each shot
  with the same Alpha Earth embedding we classify on, plus slope.
  \item Finding: biomass is a regression target on our own feature space, so it slots in as a Random
  Forest regressor on the same features, with no new pipeline.
  \item We reproduced the collection and reached the expected ballpark; it stays out of the interface
  until we decide how biomass should appear in the LULC.
\end{itemize}
\end{frame}

% ------------------------------------------------------------------
\begin{frame}{Segmenting the mining class}
\begin{itemize}
  \item The mining class is per-pixel, so it shows as scattered pixels; the ask was mining as discrete
  objects.
  \item This is \emph{not} the object-detection model intended for this: that needs imagery, masks, and
  GPU we have not wired up. As an interim that reuses the existing tile path, we vectorise the mining
  prediction in Earth Engine: de-speckle, close holes, trace polygons, drop tiny blobs.
  \item Result: clean mining polygons with a per-object area, downloadable, instead of pixel confetti,
  useful now. The learned object-detection model is the real next step.
\end{itemize}
\end{frame}

% ------------------------------------------------------------------
\begin{frame}{Water, step one: robust across bodies and years}
Two steps: first get water against non-water working and see the accuracy, then decide how to fold it
into the LULC. This is step one.
\begin{itemize}\small
  \item We hold out whole water bodies (spatial), whole years (temporal), and the combination (the
  honest worst case). The polygons already carry a water-body identity and a date.
  \item Within water bodies the task is strong and stable: the combined spatial-and-temporal hold-out
  scores about 0.98, with a small year-to-year spread, so no fluke year.
  \item So the classifier is solid enough within water bodies to build on.
\end{itemize}
\end{frame}

% ------------------------------------------------------------------
\begin{frame}{Water, step one: works anywhere, and counting fortnights}
\begin{itemize}\small
  \item The catch: the non-water examples all come from around water bodies, so the model over-calls
  water on dry land. We augment non-water with barren, built-up, and greenery pixels across seasons.
  \item The effect: non-water precision rises from about 0.72 to 0.99, so it stops painting dry land as
  water. This works-anywhere model is now deployed.
  \item We also run the model over a year of fortnights and count, per pixel, how many fortnights held
  water, so perennial bodies score high and monsoon-only ponds low. This is the seasonal-water signal the
  LULC would use, deferred to step two.
\end{itemize}
\end{frame}

% ------------------------------------------------------------------
\begin{frame}{Acacia: measuring spatial and temporal robustness}
Acacia against non-acacia is a hard species split (two similar trees). The point here is the honest
measurement, not a headline number: crowns now carry their source region, so we can hold out a whole
region and whole years at once.
\vspace{2mm}
\centering
\begin{tabular}{lr}
\toprule
What is held out & Accuracy \\
\midrule
Unseen years only (temporal) & 0.749 \\
An unseen region only (spatial) & 0.695 \\
An unseen region \emph{and} unseen years & 0.679 \\
\bottomrule
\end{tabular}
\vspace{2mm}
\begin{flushleft}\small
The ordering is what it should be: a whole region is harder than random polygons, and region-and-year
together is hardest. That gap is the value, since it shows how much a single-site, single-year acacia
number would over-promise. Reassuringly, across the two test years the accuracy is stable (spread
$\sim$0.002), so there is no fluke year: the split travels consistently, just modestly, because it is a
hard species distinction.
\end{flushleft}
\end{frame}

% ------------------------------------------------------------------
\begin{frame}{How long does Tessera take, against Alpha Earth}
Measured live over one small site, to advise which pipeline to reach for.
\vspace{2mm}
\centering
\begin{tabular}{lrr}
\toprule
Stage & Tessera & Alpha Earth \\
\midrule
Download & $\sim$29\,s / tile (151\,MB) & none, server-side \\
Sample & 31\,s & 4\,s \\
Train & $<$1\,s & $<$1\,s \\
Classify & 40\,s, point grid & 8\,s, tiles \\
\midrule
Total & $\sim$72\,s & $\sim$12\,s \\
\bottomrule
\end{tabular}
\vspace{2mm}
\begin{flushleft}\small
Alpha Earth is server-side round-trips only. Tessera pays a large per-tile download first, then samples
and classifies locally. So Alpha Earth wins on first touch; Tessera is worth it only when its local
features are needed and the download is already paid.
\end{flushleft}
\end{frame}

% ------------------------------------------------------------------
\begin{frame}{STACD provenance for every output}
STACD provenance was implemented for every classified output, following the paper's five classes.
\begin{itemize}\small
  \item Each output emits a STAC Item for the raster (box, geometry, class legend, assets) plus the
  STACD graph: a node per input dataset and algorithm, one instance per live model, rule, and merge,
  and the output instance.
  \item The class hierarchy and the ordered operations that built it are embedded as the producing
  algorithm's input set, reusing the zoo cards, so there is no new source of truth and no
  Earth-Engine run.
  \item We emit the metadata half; the paper's Airflow runtime is the next layer. The record lines up
  with the drone and bioacoustics outputs, so it can be compared across the projects.
\end{itemize}
\end{frame}

% ------------------------------------------------------------------
\begin{frame}
\centering
\vspace{1.5cm}
{\Large Thank you}
\vspace{0.5cm}


\end{frame}

\end{document}
